\section{Discussion}
\label{sec:discussion}

The central contribution of this paper is to expose a semantic mismatch between
cross-validation diagnostics and adversarial replacement robustness for
deterministic local learning rules. Leave-one-out stability is a statement about
the effect of deleting a point and evaluating at that same point; replace-one
stability concerns the effect of replacing a point by an adversarially chosen
labeled example and evaluating at any query. The paper shows that these two
notions are structurally distinct, through a systematic perturbation calculus,
an exact margin-crossing criterion, and a general separation theorem for $k$-NN
that holds for every $k \ge 1$.

\subsection{Established Results}

\begin{enumerate}[leftmargin=*,label=(\arabic*)]
    \item \textbf{Uniform decomposition:} replace-one $\le$ delete-one +
        insert-one (\Cref{prop:4.1}); the decomposition is algebraically sharp.
    \item \textbf{Margin criterion:} deterministic $k$-NN prediction changes iff
        the signed vote margin crosses zero (\Cref{prop:5.1}).
    \item \textbf{General separation:} for every odd $k \ge 3$, there exists a
        finite metric sample with $\Delta_{\mathrm{loo}}(S,i)=0$ for all $i$ and
        $\Delta_{\mathrm{rep}}(S,j,z,x,y)=1$ (\Cref{thm:odd-k-separation}).
    \item \textbf{Even-$k$ separation:} the analogous result holds for even
        $k \ge 2$ under deterministic tie-breaking
        (\Cref{thm:even-k-separation}), with a tighter dependence on the
        tie-breaking convention.
\end{enumerate}

\subsection{Stability Hierarchy}

The relationships among the stability notions studied in this paper can be
summarized concisely. At the uniform level, \Cref{prop:4.1} shows that
delete-one and insert-one together control replace-one, but neither suffices
alone. At the pointwise level, \Cref{thm:odd-k-separation} and
\Cref{thm:even-k-separation} show that LOO stability does not imply replace-one
stability, and the converse is trivially false because of the different
evaluation quantification. Pointwise delete-one (over all queries) implies LOO,
but not conversely. Standard relationships between expected and worst-case
stability (expected does not imply worst-case; worst-case implies expected)
complete the picture.

\begin{table}[t]
\centering
\caption{Established stability relationships. Each row states an implication
between stability notions, its status (theorem or counterexample), and the
relevant reference. Rows with uncertain direction or insufficient evidence
are omitted.}
\label{tab:stability-hierarchy}
\small
\begin{tabularx}{\linewidth}{lcl}
\toprule
\textbf{Relationship} & \textbf{Status} & \textbf{Reference}\\
\midrule
Uniform delete+insert
$\Rightarrow$ uniform replace-one & Theorem & \Cref{prop:4.1}\\
\addlinespace
Pointwise LOO
$\Rightarrow$ pointwise replace-one & Counterexample & \Cref{thm:odd-k-separation},
\Cref{thm:even-k-separation}\\
Pointwise replace-one
$\Rightarrow$ pointwise LOO & Counterexample & Different evaluation
quantification\\
\addlinespace
Pointwise delete-one (all queries)
$\Rightarrow$ pointwise LOO & Theorem & Definition chase; LOO is a
special case\\
Pointwise LOO
$\Rightarrow$ pointwise delete-one (all queries) & Counterexample & Different
evaluation point\\
\addlinespace
Expected replace-one
$\Rightarrow$ worst-case replace-one & Counterexample & Standard; rare
adversarial configurations\\
Worst-case replace-one
$\Rightarrow$ expected replace-one & Theorem & Supremum dominates expectation\\
\bottomrule
\end{tabularx}
\end{table}

% Hierarchy summary table: maps each pair of notions to a concise classification.
% Only includes relationships whose direction is correctly established.
\begin{table}[t]
\centering
\caption{Summary classification of stability notion pairs.}
\label{tab:hierarchy-summary}
\small
\begin{tabularx}{\linewidth}{lccc}
\toprule
\textbf{Notion $A$} & \textbf{Notion $B$} & $A \Rightarrow B$? & $B \Rightarrow A$?\\
\midrule
Uniform delete+insert & Uniform replace-one & Yes (\Cref{prop:4.1}) & Open\\
Pointwise LOO & Pointwise replace-one & No (\Cref{thm:odd-k-separation}) & No\\
Pointwise delete-one (all) & Pointwise LOO & Yes & No\\
Expected replace-one & Worst-case replace-one & No & Yes\\
\bottomrule
\end{tabularx}
\end{table}


Several relationships remain open at the uniform level, including the
comparability of uniform delete-one and uniform insert-one, and the converse of
the delete-plus-insert decomposition.

\subsection{Relationship to Other Notions of Stability}

It is useful to situate the present work relative to established lines of
inquiry that also use the term ``stability.''

\paragraph{Generalization stability (Bousquet--Elisseeff).}
Classical algorithmic stability \cite{bousquet2002stability} studies how
replacing one training example affects the expected or fresh-test loss, with
the goal of bounding generalization error. That agenda is distributional and
asymptotic. Our work is combinatorial and worst-case: we fix a finite sample
and ask whether a single perturbation changes the prediction at a specific
query. There is no contradiction between a rule being stable in the
generalization sense and having the pointwise LOO/replace-one gap we
demonstrate.

\paragraph{Consistency (Cover--Hart, Stone).}
Classical $k$-NN theory \cite{cover1967nearest,stone1977consistent} studies
whether the classifier converges to the Bayes risk as $n \to \infty$ under
distributional assumptions. Those are limit theorems about expected performance.
Our witness constructions are finite and adversarial. A rule can be consistent
and simultaneously admit finite configurations where LOO and replace-one
diverge.

\paragraph{Conformal prediction stability.}
Recent work \cite{liang2023algorithmic,pournaderi2024training} studies how
deleting or exchanging points affects prediction sets and coverage certificates.
Those analyses operate under exchangeability or training-conditional assumptions
and typically consider real-valued conformity scores rather than 0--1 loss. Our
results are complementary: they show that even in the simpler 0--1 loss setting,
the choice of perturbation operation matters before any distributional or
coverage considerations enter.

\subsection{Practical Implications}

The LOO/replace-one gap has concrete consequences for practitioners using
$k$-NN or similar local rules. Low LOO error is often interpreted as evidence of
model stability, but a model can be perfectly LOO-stable while being highly
vulnerable to a single adversarially chosen replacement. The margin-based
diagnostic provides a computationally efficient way to detect vulnerable
queries: queries with $|M_k(S,x)| \le 2$ are potentially vulnerable to a single
label-flip replacement. Practitioners should therefore not equate LOO stability
with replacement robustness, and should monitor the margin distribution rather
than relying solely on aggregate LOO error.

\subsection{Limitations}

\begin{enumerate}[leftmargin=*,label=(\arabic*)]
    \item \textbf{Beyond binary 0--1 loss.} The calculus is developed for binary
        classification with 0--1 loss. Extending it to real-valued scores,
        multiclass labels, or regression loss functions would require a more
        general notion of margin crossing.
    \item \textbf{Beyond finite metric spaces.} The constructions use finite
        metric spaces (often graph metrics). Generalizing to infinite spaces,
        continuous densities, or non-metric similarity measures is open.
    \item \textbf{Beyond $k$-NN.} The uniform calculus of \Cref{prop:4.1}
        extends to any learning rule with a bounded loss, but the explicit
        separation constructions use $k$-NN combinatorics. Finding analogous
        constructions for other local rules would strengthen the claim of
        generality.
    \item \textbf{Randomized tie-breaking.} Our analysis is for deterministic
        tie-breaking. The even-$k$ results show that randomization can close
        some gaps. Characterizing the stability hierarchy under randomized
        tie-breaking is an open problem.
    \item \textbf{Distributional connections.} The paper studies worst-case
        fixed-sample stability. Connecting the finite combinatorial guarantees
        to distributional stability notions without collapsing the distinctions
        we isolate is an important next step.
\end{enumerate}

These limitations also point to follow-up problems. Strengthening the
certificate pipeline into a formal minimality argument, characterizing which
metric-space constraints eliminate the separation, extending the perturbation
calculus to other local rules, and connecting the finite worst-case indicators
here to distributional stability notions are all natural directions.

\subsection{Open Problems}

\begin{enumerate}[leftmargin=*,label=(\arabic*)]

    \item \textbf{Tie-free even-$k$ separation.}
    Does there exist, for any even $k \ge 2$, a finite metric sample for
    deterministic $k$-NN in which the original prediction at the designated
    query is determined by a strict majority (margin $\neq 0$) and the
    LOO/replace-one separation holds? The constructions in this paper all use
    margin $0$ for even $k$, which ties the separation to tie-breaking.

    \item \textbf{Quantitative bounds.}
    For a given data distribution, what is the expected LOO/replace-one gap?
    Can the worst-case gap be bounded by a function of the data geometry,
    such as the doubling dimension or the number of duplicate-prone points?

    \item \textbf{Algorithmic mitigation.}
    Can the margin-based diagnostic be extended to an algorithm that repairs
    vulnerable queries, for instance by recommending a different $k$ or a
    specific subset of points to re-label?

    \item \textbf{Connection to adversarial robustness.}
    The replace-one perturbation is a form of training-set adversarial example.
    How does the LOO/replace-one gap relate to standard test-time adversarial
    examples? Is there a transfer between training-set and test-set
    vulnerabilities?

\end{enumerate}
